\documentclass[11pt,a4paper]{article}
\usepackage[ngerman]{babel}
\usepackage[T1]{fontenc}
\usepackage[utf8]{inputenc}
\usepackage{amsmath,amssymb,mathtools}
\usepackage[most]{tcolorbox}
%\usepackage{geometry}
\usepackage{enumitem}
\usepackage{microtype}
%\geometry{margin=2.35cm}
%\setlength{\parskip}{0.45em}
%\setlength{\parindent}{0pt}
%\setlist[enumerate]{leftmargin=*,itemsep=0.25em}
\tcbset{colback=gray!5,colframe=gray!65!black,arc=1mm,boxrule=0.5pt,breakable,enhanced}

\title{Moral und (A)theismus}
\author{Daniel Zahnd, Nicola Liebi}
\date{\today}

\begin{document}
	\maketitle
	
	\section{Einleitung}
	
	In diesem Artikel wird ein zentrales Argument gegen den evolutionären Naturalismus und für den Theismus vorgestellt. Ein Teil des Arguments basiert auf Mark D. Linvilles \textit{Argument from Evolutionary Naturalism}. Bei diesem Argument spielt die Existenz objektiver moralischer Fakten eine entscheidende Rolle. Mit objektiven moralischen Fakten sind Aussagen gemeint, die unabhängig von Meinungen, Gefühlen oder kulturellen Überzeugungen wahr sind. Ziel dieses Artikels ist es aufzuzeigen, dass ein naturalistisch-materialistisches Weltbild inkompatibel mit der Existenz objektiver moralischer Fakten ist. Das Argument wird anschliessend erweitert, um zu zeigen, dass objektive moralische Fakten auf die Existenz eines göttlichen Wesens (Gott) hindeuten. Die Argumentationsstruktur ist wie folgt:
	
	\begin{tcolorbox}[title=Die vier Prämissen des Arguments]
		\begin{enumerate}[label=(\arabic*)]
			\item Wenn der evolutionäre Naturalismus wahr ist, sind moralische Fakten Produkte natürlicher Selektion.
			\item Wenn moralische Fakten Produkte natürlicher Selektion sind, existieren keine objektiven moralischen Fakten.
			\item Es existieren objektive moralische Fakten.
			\begin{enumerate}[label=(\arabic*)]
				\item[] (1), (2), (3) $\rightarrow$ (S1)$\,$ Der evolutionäre Naturalismus ist falsch.
			\end{enumerate}
			\item Wenn kein göttliches Wesen existiert, dann ist der evolutionäre Naturalismus wahr.
			\begin{enumerate}[label=(\arabic*)]
				\item[] (S1), (4) $\rightarrow$ (S2)$\,$ Ein göttliches Wesen existiert.
			\end{enumerate}
		\end{enumerate}
	\end{tcolorbox}
	
	Aus den Prämissen (1), (2) und (3) folgt logisch zwingend, dass der evolutionäre Naturalismus falsch ist. Nimmt man Prämisse (4) hinzu, folgt daraus, dass ein göttliches Wesen existiert. Das bedeutet: Wenn alle vier Prämissen akzeptiert werden, dann folgt daraus, dass Gott existiert. Nachfolgend werden die Argumente für die einzelnen Prämissen genauer betrachtet. Ausklappbare Boxen ergänzen mathematische Details, wie die Plausibilisierung der jeweiligen Prämisse gezeigt werden kann.
	
	\begin{tcolorbox}[title=Mathematik-Box: Die logische Struktur des Arguments]
		Es seien
		\[
		\begin{aligned}
			E &:= \text{Der evolutionäre Naturalismus ist wahr},\\
			N &:= \text{Moralische Fakten sind Produkte natürlicher Selektion},\\
			M &:= \text{Objektive moralische Fakten existieren},\\
			G &:= \text{Ein göttliches Wesen existiert}.
		\end{aligned}
		\]
		Dann lauten die Prämissen formal:
		\[
		p_1:E\rightarrow N,\qquad
		p_2:N\rightarrow\neg M,\qquad
		p_3:M,\qquad
		p_4:\neg G\rightarrow E.
		\]
		Aus $p_1$, $p_2$, und $p3$ folgt  \(\neg E\). Mit \(p_4\) ergibt sich
		\[
		\neg G\rightarrow E\rightarrow N\rightarrow\neg M.
		\]
		Für eine Aussage $\neg A \rightarrow \neg B$ ist die Kontraposition $B \rightarrow A$ immer wahr. Da nun durch Prämisse (3) \(M\) angenommen wird, folgt daraus durch Kontraposition, dass \(\neg G\) nicht wahr sein kann und deshalb
		\( M \rightarrow G\) gilt. Sind die Prämissen wahr, dann ist die Schlussfolgerung $G$ also unumstösslich wahr. Nun kann aber nicht zweifelsfrei gezeigt werden, dass die Prämissen wahr sind; es kann aber dafür argumentiert werden, dass sie mit grösserer Wahrscheinlichkeit wahr sind, als dass sie falsch sein könnten. Durch Argumentation können die Prämissen plausibilisiert und damit gezeigt werden, dass daher auch die Schlussfolgerung $G$ mit grösserer Wahrscheinlichkeit wahr ist, als falsch zu sein.
	\end{tcolorbox}
	
	\section{Diskussion von Prämisse (1)}
	
	\textit{„Wenn der evolutionäre Naturalismus wahr ist, sind moralische Fakten Produkte natürlicher Selektion.“}
	
	Diese Prämisse sagt aus, dass moralische Vorstellungen im evolutionären Naturalismus als Produkte der natürlichen Selektion angesehen werden müssen. Anders ausgedrückt müssen moralische Überzeugungen im Rahmen eines evolutionär-naturalistischen Weltbildes als durch natürliche Selektion hervorgebrachte Eigenschaften verstanden werden. Das bedeutet: Moral hat keinen metaphysischen\footnote{In diesem Zusammenhang bedeutet „metaphysisch“, dass etwas über die rein physische Welt hinausgeht – also nicht letztlich physikalisch erklärbar ist.} Ursprung, sondern einen evolutionären Nutzen.
	
	Im materialistischen Naturalismus muss alles auf Materie und deren Gesetzmässigkeiten zurückgeführt werden, der sie unterliegt. Gemäss dieser Anschauung sind geistige Phänomene wie moralische Vorstellungen nur Epiphänomene\footnote{  Also Das heisst, „Begleiterscheinungen“. Kurz gesagt: Im materialistischen Naturalismus sind moralische Gedanken nicht eigenständig wirksam, sondern nur Folgeerscheinungen der materiellen Vorgänge im Körper.} materieller Prozesse. Auch psychische Ereignisse sind letztlich durch physische Ursachen determiniert.
	
	Daraus folgt: Moralische Überzeugungen entstehen aufgrund letztlich physikalischer Prozesse im Körper. Diese Prozesse sind durch die Beschaffenheit des Körpers determiniert. Dieser wiederum wird stark durch das Zusammenspiel natürlicher Selektion und zufälliger Mutationen geprägt.
	
	Es ist leicht zu sehen, inwiefern moralische Regeln Selektionsvorteile bieten können. Beispielsweise bieten die Fürsorgepflicht gegenüber Kindern oder das Verbot von Mord evolutionäre Vorteile, da diese (1) die Überlebenschancen der Gruppe erhöhen und (2) die Kooperation und soziale Stabilität stärken. So machen moralische Regeln die Gruppe insgesamt überlebensfähiger gegenüber dem Prozess der natürlichen Selektion.
	
%	\subsection*{„Greedy Reductionism“}
%	
%	Doch können moralische Überlegungen effektiv rein auf physische bzw. biologische Prozesse reduziert werden? Daniel Dennett warnt vor einer überzogenen Reduktion aller geistigen Phänomene (wie moralisches Urteil) auf biologische Prozesse. Er bringt ein Beispiel: Dass alle Kulturen Speere mit der Spitze voran werfen, beweise nicht, dass es ein entsprechendes Gen gäbe. Moralisches Verhalten müsse also nicht zwingend genetisch verankert sein.
%	
%	Die Argumentation im PDF versucht darauf zu antworten, dass bei einem materialistischen Naturalismus auch geistige Phänomene letztlich durch materielle Vorgänge determiniert sein müssten. Dies bedeutet nicht zwingend, dass es für jede einzelne moralische Überzeugung ein bestimmtes Gen gibt; die These ist vielmehr, dass moralische Vorstellungen letztlich auf natürliche Prozesse zurückgeführt werden.
	
	\begin{tcolorbox}[title=Mathematik-Box: Prämisse 1]
		Es stehe \(V\) für „Population \(\mathcal{P}\) hat Moralvorstellungen“ und \(S\) für „Population \(\mathcal{P}\) hat einen Selektionsvorteil“. Was plausiblisiert werden soll ist die Implikation $V \rightarrow S$, also dass Moralvorstellungen einer Population ihr einen Selektionsvorteil einbringen, d. h. dass Moralvorstellungen Produkte natürlicher Selektion sind. 
		
		In der sogenannten Odds-Form lautet Bayes' Theorem:
		\[
		\frac{P(S\mid V)}{P(\neg S\mid V)}
		=
		\frac{P(S)}{P(\neg S)}
		\frac{P(V\mid S)}{P(V\mid\neg S)}.
		\]
		Hierbei bedeutet eine Formulation $P(A|B)$ für zwei Sachverhalte $A$ und $B$ die „Wahrscheinlichkeit von $A$, gegeben $B$“.
		
		Die Argumentation setzt zunächst \(P(S)/P(\neg S)\approx 1\), da die Wahrscheinlichkeiten von $S$ und $\neg S$ a priori dieselben zu sein scheinen. Aufgrund der gegebenen Argumentation im Fliesstext ist anzunehmen, dass Moralvorstellungen bei Populationen mit Selektionsvorteil wahrscheinlicher sind, also
		\[
		\frac{P(V\mid S)}{P(V\mid\neg S)}> 1,
		\]
		woraus
		\[
		\frac{P(S\mid V)}{P(\neg S\mid V)}>1.
		\]
		folgt. Nimmt man fernerhin an, dass eine Population a priori mit gleicher Wahrscheinlichkeit keine Moralvorstellungen hat, als dass sie welche hat, so folgt $P(V)/P(\neg V) \approx 1$. Die Odds-Form für die zu plausibilisierende Prämisse lautet nun
		$$\frac{P(S\mid V)}{P(S\mid\neg V)} = \underbrace{\frac{P(V)}{P(\neg V)}}_{\approx \,1} \underbrace{\frac{P(S|V)}{P(S|\neg V)}}_{>\,1} > 1.$$
		
		Die natürliche Selektion bringt also gemäss einem materialistisch-naturalistischen Weltbild mit einer Wahrscheinlichkeit grösser $1/2$ moralische Fakten hervor.
	\end{tcolorbox}
	
	\section{Diskussion von Prämisse (2)}
	
	\textit{„Wenn moralische Fakten Produkte natürlicher Selektion sind, existieren keine objektiven moralischen Fakten.“}
	
	\textbf{Hier weiterfahren.}
	
	Die Grundidee der zweiten Prämisse ist, dass moralische Glaubenssätze, die durch evolutionäre Prozesse entstanden sind, epistemisch unabhängig von deren Wahrheit sind. Dies wird auch die „Unabhängigkeitsthese“ genannt. Anders ausgedrückt: Die Herkunft eines Glaubenssatzes hat nichts mit der Wahrheit zu tun.
	
	Ein Beispiel zur Illustration: Ben zieht zufällig die Zahl 73 aus einer Urne und glaubt daher, dass 73 Studenten in seiner Klasse sind. Auch wenn das zufällig stimmen sollte, ist der Glaube epistemisch nicht gerechtfertigt. Übertragen auf Moral bedeutet dies Folgendes: Wenn moralische Überzeugungen durch evolutionäre Mechanismen entstehen, heisst das nicht, dass sie auch wahr sind.
	
	Anhand eines Gedankenexperiments soll dieser Punkt illustriert werden. Stellen Sie sich zwei Paare in unterschiedlichen Umwelten vor. In einer Umgebung (\(U_A\)) ist es nützlich, Kinder zu töten (Nahrungsmangel). In einer anderen (\(U_B\)) ist es nützlich, möglichst viele Kinder zu bekommen. Die daraus entstehenden Moralvorstellungen (\(M_A\) und \(M_B\)) sind einander komplett entgegengesetzt: In \(U_A\) ist es moralisch gut, Kinder zu töten, während dies in \(U_B\) schlecht ist. Wenn beide durch natürliche Selektion entstanden sind, kann keine von ihnen objektiv wahr sein. Dies soll zeigen, dass moralische Wahrheiten im Rahmen des Naturalismus kontingent und nicht objektiv sind.
	
	\begin{tcolorbox}[title=Mathematische Box: Bayes' Theorem und Prämisse (2)]
		Mit \(M\) für „Es existieren objektive moralische Fakten“ und \(N\) für „Moralische Fakten sind Produkte des evolutionären Naturalismus“ lautet die Odds-Form:
		\[
		\frac{P(M\mid N)}{P(\neg M\mid N)}
		=
		\frac{P(M)}{P(\neg M)}
		\cdot
		\frac{P(N\mid M)}{P(N\mid\neg M)}.
		\]
		Das PDF nimmt an, dass objektive Moral zunächst plausibel ist, also
		\[
		\frac{P(M)}{P(\neg M)}\gg 1.
		\]
		Zugleich wird angenommen, dass eine evolutionäre Erklärung moralischer Fakten bei objektiver Moral wenig wahrscheinlich sei, während sie bei nichtobjektiver Moral plausibler sei:
		\[
		P(N\mid M)\ll 1,\qquad P(N\mid\neg M)\text{ ist vergleichsweise hoch}.
		\]
		Wenn der kleine Faktor \(\frac{P(N\mid M)}{P(N\mid\neg M)}\) die anfängliche Einschätzung ausreichend überwiegt, ergibt sich:
		\[
		\frac{P(M\mid N)}{P(\neg M\mid N)}<1.
		\]
		Das soll die Prämisse \(N\rightarrow\neg M\) plausibilisieren. Es ist jedoch keine rein mathematische Widerlegung objektiver Moral: Das Ergebnis hängt von den angenommenen Wahrscheinlichkeiten und damit von der inhaltlichen Argumentation ab. Die Entstehung einer Überzeugung beweist für sich allein weder ihre Wahrheit noch ihre Falschheit.
	\end{tcolorbox}
	
	\section{Diskussion von Prämisse (3)}
	
	\textit{„Es existieren objektive moralische Fakten.“}
	
	Die meisten Menschen sind intuitiv von der Objektivität moralischer Fakten überzeugt. So würden die meisten Menschen zum Beispiel Folter von unschuldigen Kindern als moralisch falsch beurteilen. Mit objektiv ist dabei gemeint, dass diese unabhängig von der Meinung von Menschen so ist. Bereits ein einziges objektiv richtiges moralisches Faktum (zum Beispiel „Folter an unschuldigen Kindern ist objektiv falsch“) genügt zur Bestätigung dieser Prämisse.
	
	Wenn man die Objektivität ablehnen würde, dann kann dies zu problematischen Konsequenzen führen. Beispielsweise könnten dann selbst Handlungen wie Mord oder Folter lediglich als subjektiv falsch betrachtet werden. Das heisst, eine Person findet es vielleicht moralisch falsch und eine andere findet es moralisch gerechtfertigt. Dadurch verlieren moralische Urteile ihre normative Kraft. Eine sinnvolle ethische Diskussion wäre kaum möglich.
	
	Zusätzliche Informationen zu dieser Prämisse finden sich im Artikel „Wie lautet das moralische Argument für Gottes Existenz“.
	
	Wenn alle drei Prämissen akzeptiert werden, folgt logisch, dass der evolutionäre Naturalismus falsch ist. Dies öffnet unter Einbezug der Prämisse 4 die Tür für theistische Erklärungen moralischer Objektivität.
	
	\section{Diskussion von Prämisse (4)}
	
	\textit{„Wenn kein Gott existiert, dann ist der evolutionäre Naturalismus wahr.“}
	
	Diese Prämisse verknüpft sozusagen den Naturalismus mit dem Atheismus. Denn ohne Gott bleibt nach der Argumentation nur der evolutionäre Naturalismus zur Erklärung von Leben und Moral. Nach dem heutigen Stand des Wissens gibt es, gegeben dass Gott nicht existiert, keine wirklichen Alternativen zum evolutionären Naturalismus. Deswegen stimmen Alex Rosenberg – ein Naturalist – sowie Alvin Plantinga – ein Theist – darin überein, dass die biologische Evolutionstheorie als wesentlicher Bestandteil des evolutionären Naturalismus das „only game in town“ sei.
	
	Die Aussage „Wenn kein Gott existiert, dann ist der evolutionäre Naturalismus wahr“ lässt sich logisch als „Wenn A, dann B“ darstellen, wobei \(A\): „Kein Gott existiert“ und \(B\): „Der evolutionäre Naturalismus ist wahr“. Nach den Regeln der Logik ergibt sich daraus die gültige Kontraposition: „Wenn der evolutionäre Naturalismus nicht wahr ist, dann existiert ein Gott“.
	
	\begin{tcolorbox}[title=Mathematische Box: Bayes' Theorem und Prämisse (4)]
		Das PDF untersucht die Plausibilität von \(\neg G\rightarrow E\), wobei \(G\) für „Gott existiert“ und \(E\) für „Der evolutionäre Naturalismus ist wahr“ steht. Bayes' Theorem in Odds-Form ergibt:
		\[
		\frac{P(G\mid E)}{P(\neg G\mid E)}
		=
		\frac{P(G)}{P(\neg G)}
		\cdot
		\frac{P(E\mid G)}{P(E\mid\neg G)}.
		\]
		Umgestellt:
		\[
		\frac{P(E\mid G)}{P(E\mid\neg G)}
		=
		\frac{P(\neg G)}{P(G)}
		\cdot
		\frac{P(G\mid E)}{P(\neg G\mid E)}.
		\]
		Das PDF nimmt für die Diskussion an, dass die Ausgangseinschätzung die atheistische Hypothese begünstigt, also \(P(\neg G)/P(G)\gg 1\). Ferner wird für den materialistisch verstandenen evolutionären Naturalismus \(P(G\mid E)\approx 0\) und \(P(\neg G\mid E)\approx 1\) angenommen. Daraus wird die qualitative Abschätzung
		\[
		\frac{P(E\mid G)}{P(E\mid\neg G)}\approx 0
		\]
		abgeleitet. In der Logik des PDFs soll dies die Annahme stützen, dass \(E\) unter \(\neg G\) viel wahrscheinlicher ist als unter \(G\).
		
		Diese Rechnung liefert keine objektiv gemessenen Wahrscheinlichkeiten. Sie hängt davon ab, wie die Hypothesen definiert und die bedingten Wahrscheinlichkeiten eingeschätzt werden. Insbesondere ist die Behauptung, ohne Gott bleibe nur der evolutionäre Naturalismus, eine zusätzliche philosophische Annahme und kein Ergebnis des Bayes-Theorems allein.
	\end{tcolorbox}
	
	\section{Schlussfolgerung}
	
	Die Prämissen 1 bis 3 sollen zeigen, dass ein evolutionär-naturalistisches Weltbild nicht in der Lage ist, objektive moralische Fakten zu erklären. Daher erscheint es als weniger plausibel, wenn man an die Objektivität der Moral glaubt. Durch die Erweiterung des Arguments mit der Prämisse 4 wird aufgezeigt, dass ein theistisches Weltbild diese Lücke überzeugend füllen kann. Somit ergibt sich ein Argument gegen den evolutionären Naturalismus und für die Existenz Gottes, sofern objektive Moral angenommen wird und die übrigen Prämissen überzeugen.
	
	Dabei ist zu unterscheiden zwischen der formalen Gültigkeit des Schlusses und der Wahrheit seiner Prämissen. Die logische Struktur zeigt, was aus den Prämissen folgt; die probabilistischen Boxen stellen dar, wie das PDF einzelne Prämissen plausibilisieren möchte. Die verwendeten Wahrscheinlichkeiten sind jedoch argumentative Einschätzungen und keine empirisch bestimmten Werte.
	
	\begin{thebibliography}{9}
		\bibitem{dennett}
		Dennett, D. C. (1995).
		\textit{Darwin's Dangerous Idea: Evolution and the Meanings of Life}.
		Simon \& Schuster.
		
		\bibitem{joyce}
		Joyce, R. (2006).
		\textit{The Evolution of Morality}.
		MIT Press.
		
		\bibitem{linville}
		Linville, M. D. (2009).
		The Moral Argument.
		In: W. L. Craig und J. P. Moreland (Hrsg.),
		\textit{The Blackwell Companion to Natural Theology}, S. 391--448.
		
		\bibitem{sober}
		Sober, E. (1994).
		\textit{From a Biological Point of View}.
		Cambridge University Press.
	\end{thebibliography}
	
	\paragraph{Über den Autor}
	Daniel Zahnd, Jahrgang 1995, ist seit 2020 mit Regula verheiratet. Sie leben gemeinsam mit ihren
	zwei Söhnen in Rüschegg Heubach.
	
	Mit rund 18 Jahren traf Daniel nach einem tragischen Krankheitsfall in der Familie die
	Entscheidung, Jesus Christus nachzufolgen. Das war ein Wendepunkt, der sein Leben
	prägte. Seither setzt er sich intensiv mit der Glaubwürdigkeit der Bibel auseinander. Es ist ihm ein Anliegen aufzuzeigen, dass Wissenschaft und Glaube an den Gott der Bibel beste Freunde sind.
	
	Beruflich absolvierte er eine Lehre als Automatiker und studierte anschließend Physik,
	Mathematik und Astronomie an der Universität Bern, wo er sein Studium 2022 mit einem Bachelor in Physik und 2024 mit einem Master in Physik abschloss. Im Rahmen des Studiums besuchte er auch eine Philosophievorlesung zu
	Gottesbeweisen, die er mit einer Arbeit über das moralische Argument abschloss. Aktuell
	(2024-2028) promoviert er im Bereich Ingenieurwissenschaften am Eidgenössischen Institut
	für Metrologie und der Universität Bern.
	
	\begin{figure}[h]
		\centering
		\includegraphics[width=3cm]{zada_02_co.jpg}
	\end{figure}
	
\end{document}